% =====================================================================
\section{Problem Formulation}
\label{sec:method}
% =====================================================================

Consider $N$ observations $\bx_i \in \Reals^d$ and a candidate geometric
model~$\btheta$ with $n_\theta$ degrees of freedom.  Noise-free observations
lie on the model manifold $\cM_{\btheta} = \{\bx : \bg(\bx, \btheta) =
\bzero\}$ of codimension~$d_g$, and a minimal sample has
$m = \lceil n_\theta/d_g \rceil$ observations.  As in~\cite{Pritts2026},
$\br_i \in \Reals^{d_g}$ is the whitened residual of observation~$i$ and
$\delta_i = \|\br_i\|$ its distance to the model.

\subsection{The extreme-value profile-marginal score}
\label{sec:evpm}

We score~$\btheta$ by a \gaussunif mixture in which both components are
fitted to the data.  A latent label $z_i \in \{0,1\}$ marks observation~$i$
as inlier or outlier.  Grouping the labels into a partition, with inliers
$\cI$, outliers $\cO = \{1,\ldots,N\}\setminus\cI$, $n_I = |\cI|$ and
$n_O = |\cO|$, the partition-conditioned likelihood is
\begin{equation}
	p(\br \mid \btheta, \cI, \sigma)
	= \prod_{i\in\cI} \pin(\br_i;\sigma) \cdot \pout(\cO)^{\,n_O},
	\label{eq:partition-likelihood}
\end{equation}
with the Gaussian inlier density
$\pin(\br_i;\sigma) = (2\pi\sigma^2)^{-d_g/2}e^{-\|\br_i\|^2/2\sigma^2}$ of
unknown inlier scale~$\sigma$, and the uniform outlier density
\begin{equation}
	\pout(\cO) = \bigl(2\,\ahat(\cO)\bigr)^{-d_g}
	\label{eq:pout}
\end{equation}
on $[-\ahat, \ahat]^{d_g}$, whose half-width $\ahat(\cO)$ is estimated from
the outliers~$\cO$ themselves by extreme value theory (\cref{sec:outlier}).
The inlier scale is integrated out for each partition and the partition is
then chosen:
\begin{align}
	S(\btheta, \cI) & = \log\int_0^\infty
	p(\br \mid \btheta, \cI, \sigma)\,\pi(\sigma^2)\,d\sigma^2,
	\label{eq:evpm}                                                  \\
	S^*(\btheta)    & = \max_{\cI} S(\btheta, \cI),
	\label{eq:evpm-profile}
\end{align}
with the conjugate prior $\pi(\sigma^2) = \IG(\alpha_0, \beta_0)$ or its
non-informative Jeffreys limit.  $S^*$ is our score, the
\emph{extreme-value profile-marginal score}; every method in
\cref{sec:experiments} is compared against it.  Neither component carries a length set by
hand: the inlier scale is marginalized, and the outlier support is
estimated from the partition it describes.  The label prior is
$\gamma = \tfrac12$, which makes the prior of every partition the constant
$2^{-N}$.

\subsection{The inlier component: marginalizing the scale}
\label{sec:inlier}

The inlier component is that of the profile-marginal
score~\cite{Pritts2026}, and the integral \cref{eq:evpm} has its closed
form
\begin{equation}
	\begin{aligned}
		S(\btheta, \cI)
		 & = \log\Gamma\Bigl(\frac{\nu}{2} + \alpha_0\Bigr)
		- \Bigl(\frac{\nu}{2} + \alpha_0\Bigr)
		\log\Bigl(\frac{\RSS_\cI}{2} + \beta_0\Bigr) \\
		 & \quad - \frac{\nu}{2}\log(2\pi)
		- n_O\,d_g\log\bigl(2\,\ahat(\cO)\bigr),
	\end{aligned}
	\label{eq:evpm-closed}
\end{equation}
where $\RSS_\cI = \sum_{i\in\cI}\delta_i^2$ and $\nu = n_I d_g - n_\theta$
counts the residual coordinates left after the $n_\theta$ consumed by
fitting~$\btheta$.  The Jeffreys prior $\pi(\sigma^2) \propto 1/\sigma^2$ is
the limit $(\alpha_0, \beta_0) \to (0,0)$ and requires $\nu > 0$, i.e.\ at
least one inlier beyond the minimal sample.  With a fixed half-width~$a$ in
place of $\ahat(\cO)$, \cref{eq:evpm-closed} is the scale-marginalized
log-likelihood of~\cite{Pritts2026}, and its maximum over partitions the
profile-marginal score.

As there, we maximize over the prefixes of the distance-sorted
observations, so one sort and one linear sweep evaluate
$S^*$.  Exchanging an inlier and an outlier that both lie
below the threshold~$u$ of \cref{sec:outlier} leaves the exceedances, and
hence $\ahat$, unchanged, so the prefix argument of~\cite{Pritts2026}
applies to every such exchange.

\subsection{The outlier component: the support from the upper tail}
\label{sec:outlier}

Under \cref{eq:pout}, every coordinate of an outlier residual is uniform on
$[-a, a]$, so the pooled absolute coordinates of the complement,
\begin{equation}
	\cZ(\cO) = \bigl\{\, |r_{ij}| : i \in \cO,\ j = 1,\ldots,d_g \,\bigr\},
	\label{eq:pooled}
\end{equation}
are $n_O d_g$ draws from a distribution on $[0, a]$ whose upper endpoint is
the half-width.  Estimating the outlier support is thus estimating the
endpoint of a distribution from its upper tail, which is what extreme value
theory does.

For a threshold~$u$, the exceedances $y = z - u$ of the values $z > u$
follow, for almost any distribution and $u$ high enough, a generalized
Pareto distribution~\cite{Pickands1975,Balkema1974}
\begin{equation}
	P(Y \le y) = 1 - \Bigl(1 + \frac{\xi\,y}{\omega}\Bigr)^{-1/\xi},
	\label{eq:gpd}
\end{equation}
with scale $\omega > 0$ and shape~$\xi$.  A bounded support, as for outlier
residuals inside the image, has $\xi < 0$ and the endpoint
$u + \omega/|\xi|$; a uniform tail has $\xi = -1$.  We set the threshold to
the quantile of the complement at level~$\eta$, $u = Q_\eta(\cZ(\cO))$, fit
$(\xi, \omega)$ to the exceedances by probability-weighted
moments~\cite{Hosking1987}, and take the endpoint as the half-width:
\begin{equation}
	\ahat(\cO) = u + \frac{\hat\omega}{|\hat\xi|}.
	\label{eq:endpoint}
\end{equation}
The quantile level~$\eta$, for example $0.9$, is the only choice the score
makes.

\subsection{Properties}
\label{sec:properties}

\begin{proposition}[No length parameter]
	\label{prop:invariance}
	Under the Jeffreys prior, the partition that maximizes
	$S(\btheta, \cdot)$, and the ranking of candidate models
	by $S^*$, are invariant to rescaling all residuals by any
	$c > 0$.
\end{proposition}
\begin{proof}
	Rescaling multiplies $\RSS_\cI$ by $c^2$.  The quantile~$u$, the
	exceedances and the fitted scale $\hat\omega$ scale with~$c$ and the
	shape $\hat\xi$ does not, so $\ahat(\cO)$ scales with~$c$.  In the Jeffreys
	limit of \cref{eq:evpm-closed}, the second term changes by $-\nu\log c$
	and the last by $-n_O d_g\log c$, a total of
	$-(n_I d_g - n_\theta + n_O d_g)\log c = -(N d_g - n_\theta)\log c$,
	the same for every partition and every candidate.
\end{proof}
With an informative prior, $\beta_0$ has units of length squared; it is
invariant in the same sense when, as in~\cite{Pritts2026}, it is fitted by
empirical Bayes to residuals in the same units.  The profile-marginal score
with a fixed~$a$ has no such property: its outlier cost does not change
with~$c$, so the shift $-\nu\log c$ differs between partitions.

\paragraph{The implicit threshold.}
The boundary of the optimal partition is the implicit
threshold~\cite{Pritts2026}
\begin{equation}
	\tau^* = \hat s\,\sqrt{2\,d_g\log\bigl(\ahat\sqrt{2/\pi}/\hat s\bigr)},
	\label{eq:implicit-threshold}
\end{equation}
where $\hat s^2 = (2\beta_0 + \RSS_{\cI^*})/(2\alpha_0 + \nu)$ is the
variance estimate of the optimal partition~$\cI^*$.  With a fixed~$a$,
$\tau^*/\hat s$ depends on the units of the residuals; with $\ahat$, it
depends only on the dimensionless ratio $\ahat/\hat s$, which the data
determine.

\paragraph{Insensitivity to the threshold quantile.}
The quantile level~$\eta$ replaces the half-width~$a$ as the score's one
choice, but the two are not alike.  The half-width is a length that sets
the outlier cost directly.  The quantile level is dimensionless, and it does
not change what is estimated: if the exceedances over~$u_0$ are generalized
Pareto with parameters $(\xi, \omega_0)$, then those over any $u > u_0$ are
generalized Pareto with $(\xi, \omega_0 + \xi(u - u_0))$~\cite{Coles2001},
and the endpoint $u + (\omega_0 + \xi(u - u_0))/|\xi| = u_0 + \omega_0/|\xi|$
does not depend on~$u$.  Raising~$\eta$ only trades bias for variance:
fewer exceedances, a noisier fit, the same endpoint.  \Cref{sec:synthetic}
confirms that the selected partition is unchanged across
$\eta \in [0.8, 0.95]$.

\paragraph{The tail shape as a diagnostic.}
The fitted shape~$\hat\xi$ tests the outlier component at no extra cost.  A
value near~$-1$ means that the far outliers are uniform up to an edge, as
\cref{eq:pout} assumes; a clearly larger value, such as $-\tfrac12$ for a
density that falls linearly to zero, signals that the outliers thin out and
that the endpoint overstates the outlier density near the model.

\paragraph{Why the upper tail and not the minima.}
\label{sec:why-maxima}
The minima of the complement are the extremes nearest the model, and an
estimate built from them, for example from the minima of random blocks of
the complement, seems attractive: it measures the outlier density where the
inlier/outlier decision is made.  It fails.  For a short prefix the
complement still contains the true inliers, its minima are inlier
distances, the estimated half-width collapses, and the outlier cost
$-n_O d_g\log(2\ahat)$ rewards labelling every point an outlier; the score
selects its smallest feasible partition (\cref{sec:synthetic}).  The upper
tail is formed by the gross outliers far from the model, which do not move
as the sweep passes through the inliers, so the estimate cannot adapt to
the points it classifies.

\subsection{Computation}
\label{sec:computation}

The sweep visits the prefixes in increasing size, so the complement only
loses its smallest values.  The exceedances over a quantile of the
complement change by at most a few values per step, and the
probability-weighted moments, which are rank-weighted sums, can be updated
with them.  \TODO{state and verify the incremental update and its cost;
alternatively fit the tail once per model to the top fraction of all
residual coordinates, which are outliers for every sensible prefix, and
check that this matches the per-prefix fit}
